\chapter{Conclusiones}
\label{chap:conclusiones}

\noindent Se ha demostrado que es factible desarrollar un modelo basado en l├│gica difusa para clasificar de manera fiable el comportamiento sedentario semanal. El modelo utiliza exclusivamente datos biom├®tricos multivariados recopilados en condiciones de vida libre por dispositivos de consumo. El sistema de inferencia difusa propuesto respondi├│ afirmativamente a la pregunta de investigaci├│n al alcanzar una alta concordancia (F1-Score = 0.840) con una clasificaci├│n objetiva y emp├¡rica derivada de un an├ílisis de conglomerados, validando as├¡ su arquitectura y su capacidad para representar patrones de comportamiento reales.

La principal contribuci├│n cient├¡fica de esta tesis es la demostraci├│n emp├¡rica del poder de los sistemas expertos para modelar la sinergia entre biomarcadores. Se identific├│ una paradoja metodol├│gica donde variables individualmente d├®biles para discriminar el comportamiento (como la Variabilidad de la Frecuencia Card├¡aca) resultaron ser indispensables para el rendimiento ├│ptimo del modelo multivariado. Este hallazgo prueba que la fortaleza del sistema no reside en la suma de sus partes, sino en su capacidad para interpretar las interacciones complejas entre la actividad f├¡sica y las respuestas cardiovasculares, una capacidad que los an├ílisis estad├¡sticos univariados no logran capturar.

Metodológicamente, este trabajo aporta una estrategia de validación robusta para sistemas de clasificación en salud cuando no se dispone de un estándar de oro clínico, basada en la convergencia de un enfoque empírico (descubrimiento de patrones no supervisado) y uno experto (reglas basadas en conocimiento). Asimismo, la ingeniería de las variables Actividad Relativa y Superávit Calórico Basal se presenta como una solución eficaz para la normalización fisiológica de datos de dispositivos portátiles, lo que permite comparaciones más equitativas entre individuos.

En conclusi├│n, esta investigaci├│n no solo entrega un modelo de clasificaci├│n interpretable y de alto rendimiento, sino que tambi├®n establece un precedente metodol├│gico para el an├ílisis de datos de salud en condiciones de vida libre. El sistema desarrollado tiene el potencial de servir como una herramienta de cribado para la identificaci├│n temprana de patrones de comportamiento sedentario de riesgo, sentando las bases para futuras investigaciones orientadas hacia la personalizaci├│n de modelos y la prevenci├│n en la salud digital.
